\documentclass[a4paper,oneside]{report}
\usepackage[utf8]{inputenc}
\usepackage{graphicx}
\usepackage{setspace}
\usepackage{geometry}
%\usepackage[showframe]{geometry}
\linespread{1.6}
\begin{document}



\begin{titlepage}
\begin{center}

\begin{Large}
	\textbf{POLITECNICO DI MILANO}\\
	Facoltà di Ingegneria\\
	Scuola di Ingegneria dell'Informazione
\end{Large}

\vspace{1.5cm}
\includegraphics[width=3cm]{imgs/poli.eps}

\vspace{1.5cm}
\begin{Large}
\textbf{
DESIGN AND DEVELOPMENT OF AN ASYNCHRONOUS DATA ACCESS LAYER FOR PERVASIVE NETWORKS
}
\end{Large}

\vfill

\begin{large}
\begin{onehalfspace}

	\begin{flushleft}
	\begin{tabbing}
	Relatore: \hspace{8pt} \= Ch.mo Prof. Ing. Fabio Alberto Schreiber\\
	Correlatore: \> Ing. Emanuele Panigati
	\end{tabbing}
	\end{flushleft}

	\vspace{1cm}
	\begin{flushright}
	\begin{tabbing}
	\hspace{300pt}
	\= Tesi di laurea di:\\
	\> GUIDO ROTA \\
	\> Matr. 819938\\
	\end{tabbing}
	\end{flushright}
	
\end{onehalfspace}
\end{large}

\vspace{2cm}
Anno Accademico 2014 - 2015

\end{center}
\end{titlepage}


\chapter*{Abstract}
Computing devices permeate every aspect of our life. PCs, tablets, smartphones, identification badges, credit cards, smart-watches, wearable gadgets, traffic cameras, digital fitness bands, and personal medical devices like pacemakers or insulin injectors are only a few of the tools that we use every day, more or less consciously, to produce and consume information.

As the number of devices and services that surround ourselves increases, so does the level of mutual cooperation that we expect from them. We know that our smartphone will automatically show us the weather for our current location, sometimes even for our hometown when we are away (How does it know where I live? Did I ever tell it?). Navigation apps guide us through different itineraries at different times of the day, depending on current and expected traffic conditions. Outbreaks of the most common viral diseases can be traced and monitored by analysing what people are searching on the web and correlating that data with other sources like hospital records.

We have grown so accustomed to the tight level of integration between different services and information sources, that behaviours similar to those presented above are nowadays expected. User requirements have heightened, and products can fail to gain traction if they don’t exploit the data that is available around us in new and innovative ways; different computing devices and services must discover themselves and make mutual use of the information that they produce or consume, meshing together in what is called a Pervasive System.

Firstly envisioned by Mark Weiser, Pervasive Systems are connected networks of independent and heterogeneous devices, whose ultimate goal is to assist people in a way that is effectively invisible to the final user. They are the result of a post-pc era, where scores of computing gadgets are disseminated in our surroundings, enhancing our capabilities to sense the world and providing ubiquitous access to information.

From a software and hardware perspective, a Pervasive Systems is a rich and varied environment, hosting myriads of different network protocols, data formats, and incompatible cpu architectures. Tapping the ever increasing stream of data produced in such a heterogeneous context, and using it to build advanced and connected products, can easily become a daunting task. Since Weiser’s seminal paper, several endeavours have explored different techniques for simplifying the task of building and designing Pervasive Systems, many of which stemmed from the research on Wireless Sensor Networks (WSN).


\chapter*{Dedication}
To mum and dad

\chapter*{Declaration}
I declare that..

\chapter*{Acknowledgements}
I want to thank...

\tableofcontents

\appendix


\end{document}
